\documentclass[a4paper,12pt]{article}
\usepackage[T2A]{fontenc}
\usepackage[utf8]{inputenc}
\usepackage[russian]{babel}
\usepackage{amsmath,amssymb,icomma}
\usepackage{geometry,booktabs,tabularx,enumitem,parskip,fancyhdr,setspace,xcolor,tikz,hyperref}
\usetikzlibrary{arrows.meta,positioning,shapes.geometric,calc}
\geometry{left=22mm,right=22mm,top=23mm,bottom=22mm,headheight=15pt}
\definecolor{accent}{HTML}{008080}
\definecolor{darkaccent}{HTML}{004D4D}
\definecolor{textgray}{HTML}{333333}
\hypersetup{hidelinks,pdftitle={Движение по окружности и тяготение},pdfauthor={Лёвин Артём Александрович}}
\onehalfspacing
\setlength{\parindent}{0pt}
\setlength{\parskip}{5pt}
\setlist{nosep,leftmargin=*}
\pagestyle{fancy}\fancyhf{}
\fancyhead[L]{\small Движение по окружности и тяготение}
\fancyhead[R]{\small София Лузянина}
\fancyfoot[C]{\thepage}
\renewcommand{\headrulewidth}{0.4pt}
\newcommand{\key}[1]{\textcolor{darkaccent}{\textbf{#1}}}
\newcommand{\unit}[1]{\,\text{#1}}
\newcommand{\ex}[1]{\par\medskip\textbf{#1}\par}
\newcommand{\G}{G}
\begin{document}
\hypersetup{pageanchor=false}
\begin{titlepage}
\centering\color{textgray}
\vspace*{22mm}
{\Large\color{darkaccent} ФИЗИКА\par}
\vspace{13mm}
{\Huge Чек-лист\par}
\vspace{9mm}
{\LARGE\color{accent} Движение по окружности,\\ всемирное тяготение\\ и движение спутников\par}
\vspace{12mm}
{\large Кинематика и динамика кругового движения\par}
\vfill
{\large София Лузянина\par}
\vspace{8mm}
{\large 25 ноября 2026 г.\par}
{\small Занятие от 24 ноября 2026 г.\par}
\vspace{12mm}
{\small Лёвин Артём Александрович\\ эксклюзивно для Софии Лузяниной\par}
\vspace{16mm}
\begin{tikzpicture}
\draw[accent,line width=.9pt] (-6,0) .. controls (-3.8,.85) and (-2,-.85) .. (0,0) .. controls (2,.85) and (3.8,-.85) .. (6,0);
\fill[darkaccent] (0,0) circle (1.4pt);
\end{tikzpicture}
\vspace{2mm}
\end{titlepage}
\setcounter{page}{1}
\hypersetup{pageanchor=true}
\section{Равномерное движение по окружности}
При равномерном движении по окружности модуль скорости постоянен, направление скорости непрерывно меняется. Пусть $r$ --- радиус окружности, $N$ --- число оборотов за время $t$.
\begin{center}
\begin{tikzpicture}[>=Stealth,scale=.85]
\draw[accent,thick] (0,0) circle(2);
\draw[gray] (0,0)--(2,0) node[midway,below]{$r$};
\draw[gray] (0,0)--(50:2);
\draw[->,darkaccent] (0.65,0) arc(0:50:.65);
\node at (25:1.02) {$\varphi$};
\fill (0,0) circle(1.4pt) node[below left]{$O$};
\fill[accent] (50:2) circle(2.4pt) node[above right]{$A$};
\draw[->,very thick,accent] (50:2)--++(140:1.65) node[above]{$\vec v$};
\draw[->,thick,darkaccent] (50:2)--++(230:1.23) node[below right]{$\vec a_{\text{ц}}$};
\end{tikzpicture}
\end{center}
\begin{center}
\renewcommand{\arraystretch}{1.25}
\begin{tabularx}{\linewidth}{@{}l l X@{}}
\toprule Величина & Формула & Смысл и единица СИ\\\midrule
Период & $T=\dfrac{t}{N}$ & Время одного оборота, с\\
Частота & $\nu=\dfrac{N}{t}=\dfrac1T$ & Число оборотов за секунду, Гц\\
Линейная скорость & $v=\dfrac{2\pi r}{T}=2\pi r\nu$ & Длина дуги за единицу времени, м/с\\
Угловая скорость & $\omega=\dfrac{2\pi}{T}=2\pi\nu$ & Угол за единицу времени, рад/с\\
\bottomrule
\end{tabularx}
\end{center}
За $N$ полных оборотов пройдены путь $s=2\pi rN$ и угол $\varphi=2\pi N$ (в радианах). Поэтому $v=s/t$, $\omega=\varphi/t$.
\[\boxed{v=\omega r}\]
\key{Проверка:} за одинаковое время точка на большем радиусе одного диска проходит больший путь, хотя угловая скорость одинакова.
\clearpage
\section{Центростремительное ускорение}
Вектор скорости направлен по касательной к окружности, вектор центростремительного ускорения --- к её центру. При равномерном движении это полное ускорение.
\[
 a_{\text{ц}}=\frac{v^2}{r}=\omega^2r
 =\frac{4\pi^2r}{T^2}=4\pi^2r\nu^2.
\]
\ex{Пример 1. Изменение скорости}
Если при неизменном радиусе скорость выросла втрое, то
\[\frac{a_2}{a_1}=\frac{v_2^2/r}{v_1^2/r}=\left(\frac{3v_1}{v_1}\right)^2=9.\]
Ускорение увеличилось в 9 раз. При неизменной скорости $a_{\text{ц}}\sim 1/r$, при неизменной угловой скорости $a_{\text{ц}}\sim r$.
\section{Шестерни и соприкасающиеся колёса}
При качении без проскальзывания модули линейных скоростей в точке контакта одинаковы. Для двух колёс с неподвижными осями:
\[v_1=v_2,\qquad \omega_1r_1=\omega_2r_2,\qquad \nu_1r_1=\nu_2r_2.\]
\begin{center}
\begin{tikzpicture}[>=Stealth,scale=.8]
\draw[accent,thick] (0,0) circle(1.55);\draw[accent,thick] (2.4,0) circle(.85);
\draw[gray] (0,0)--(1.55,0) node[midway,below]{$r_1$};
\draw[gray] (2.4,0)--(1.55,0) node[midway,below]{$r_2$};
\fill (0,0) circle(1.5pt);\fill (2.4,0) circle(1.5pt);\fill[darkaccent] (1.55,0) circle(2pt);
\draw[->,thick,darkaccent] (1.55,0)--(1.55,1.2) node[above]{$\vec v$};
\draw[->,thick] (65:1.8) arc(65:145:1.8);
\draw[->,thick] ([shift={(2.4,0)}]120:1.08) arc(120:35:1.08);
\end{tikzpicture}
\end{center}
\ex{Пример 2. Передача вращения}
Большое колесо радиусом $20\unit{см}$ совершает $20$ оборотов за $10\unit{с}$; малое имеет радиус $10\unit{см}$. Тогда
\[\nu_1=2\unit{Гц},\qquad \nu_2=\nu_1\frac{r_1}{r_2}=4\unit{Гц}.\]
Внешние шестерни вращаются в противоположных направлениях; меньшее колесо совершает больше оборотов за то же время.
\clearpage
\section{Закон всемирного тяготения}
Два тела притягиваются с равными по модулю и противоположными по направлению силами. Для материальных точек или сферически симметричных неперекрывающихся тел:
\[\boxed{F=G\frac{Mm}{r^2}},\qquad G\approx 6{,}7\cdot10^{-11}\unit{Н}\cdot\text{м}^2/\text{кг}^2.\]
Здесь $M,m$ --- массы тел, $r$ --- расстояние между их центрами. По второму закону Ньютона ускорение тела равно $F/m$. У более массивного тела ускорение меньше, хотя силы притяжения равны по модулю.
\ex{Пример 3. Метод отношений}
Если расстояние между центрами увеличилось в 3 раза при постоянных массах, то
\[\frac{F_2}{F_1}=\frac{GMm/(3r)^2}{GMm/r^2}=\frac19.\]
Сила уменьшилась в 9 раз. При сравнении двух состояний удобно делить формулы друг на друга, чтобы сократить общие величины.
\section{Ускорение свободного падения}
Обозначим радиус планеты через $R_0$, её массу через $M$. У поверхности сферической планеты (без учёта её вращения):
\[mg_0=\frac{GMm}{R_0^2}\quad\Longrightarrow\quad \boxed{g_0=\frac{GM}{R_0^2}}.\]
Масса пробного тела сокращается, так как $m>0$.
\ex{Пример 4. Сравнение двух планет}
Масса второй планеты составляет $1/16$ массы первой, а радиус вдвое меньше. Тогда
\[\frac{g_2}{g_1}=\frac{M_2}{M_1}\left(\frac{R_1}{R_2}\right)^2=\frac1{16}\cdot 2^2=\frac14.\]
Ускорение свободного падения на второй планете в 4 раза меньше. Численные значения масс и радиусов для такого сравнения не нужны.
\key{Важно:} $G$ --- универсальная постоянная, $g_0$ --- ускорение, зависящее от массы и радиуса планеты. На высоте $h$ ускорение свободного падения также существует и обозначается $g(h)$.
\clearpage
\section{Круговая орбита спутника}
Радиус орбиты отсчитывается от центра планеты. При высоте $h$ над её поверхностью
\[\boxed{r=R_0+h}.\]
\begin{center}
\begin{tikzpicture}[>=Stealth,scale=.88]
\draw[accent,thick,fill=accent!5] (0,0) circle(1.35);
\draw[darkaccent,dashed,thick] (0,0) circle(2.4);
\fill (0,0) circle(1.5pt) node[below left]{$O$};
\draw[gray] (0,0)--(1.35,0) node[midway,below]{$R_0$};
\draw[<->,gray] (1.35,0)--(2.4,0) node[midway,below]{$h$};
\fill[darkaccent] (55:2.4) circle(2.8pt) node[above right]{спутник};
\draw[->,accent,very thick] (55:2.4)--++(145:1.4) node[above]{$\vec v$};
\draw[->,darkaccent,thick] (55:2.4)--++(235:1.3) node[right]{$\vec F,\ \vec a_{\text{ц}}$};
\end{tikzpicture}
\end{center}
\key{Главная идея.} В модели круговой орбиты спутника с массой $m\ll M$ сила тяготения обеспечивает всё центростремительное ускорение:
\[F=ma_{\text{ц}},\qquad \frac{GMm}{r^2}=m\frac{v^2}{r}.\]
Сокращая положительные величины $m$ и $r$, получаем:
\[v^2=\frac{GM}{r},\qquad \boxed{v=\sqrt{\frac{GM}{R_0+h}}}.\]
Используем $v=\omega r=2\pi r\nu=2\pi r/T$:
\[\omega=\sqrt{\frac{GM}{r^3}},\qquad \nu=\frac1{2\pi}\sqrt{\frac{GM}{r^3}},\qquad T=2\pi\sqrt{\frac{r^3}{GM}}.\]
\key{Условия:} орбита круговая; сопротивлением среды, другими телами и движением массивной планеты пренебрегаем. Для эллиптической орбиты эти формулы скорости и периода через текущий радиус применять нельзя.
\clearpage
\section{Как связаны параметры орбиты}
\begin{center}
\renewcommand{\arraystretch}{1.35}
\begin{tabularx}{\linewidth}{@{}l l X@{}}
\toprule Величина & Формула & При постоянной массе $M$\\\midrule
Сила притяжения & $F=GMm/r^2$ & $F\sim r^{-2}$ (при постоянной $m$)\\
Ускорение & $a_{\text{ц}}=GM/r^2$ & $a_{\text{ц}}\sim r^{-2}$\\
Линейная скорость & $v=\sqrt{GM/r}$ & $v\sim r^{-1/2}$\\
Угловая скорость & $\omega=\sqrt{GM/r^3}$ & $\omega\sim r^{-3/2}$\\
Частота & $\nu=\omega/(2\pi)$ & $\nu\sim r^{-3/2}$\\
Период & $T=2\pi\sqrt{r^3/(GM)}$ & $T\sim r^{3/2}$\\\bottomrule
\end{tabularx}
\end{center}
\ex{Пример 5. Изменение орбиты}
Радиус круговой орбиты увеличился в 4 раза. Тогда скорость уменьшилась в 2 раза, ускорение --- в 16 раз, а период увеличился в 8 раз:
\[\frac{v_2}{v_1}=\frac1{\sqrt4}=\frac12,\quad \frac{a_2}{a_1}=\frac1{4^2}=\frac1{16},\quad \frac{T_2}{T_1}=\sqrt{4^3}=8.\]
\section{Низкая орбита и плотность планеты}
Средняя плотность сферической планеты:
\[\rho=\frac{M}{V}=\frac{3M}{4\pi R_0^3},\qquad M=\frac43\pi R_0^3\rho.\]
Если $h\ll R_0$, то радиус низкой орбиты $r\approx R_0$. В идеализированном пределе:
\[T\approx 2\pi\sqrt{\frac{R_0^3}{G\cdot\frac43\pi R_0^3\rho}}=\boxed{\sqrt{\frac{3\pi}{G\rho}}}.\]
\ex{Пример 6. Сравнение Марса и Земли}
В задаче плотности Марса и Земли считаются приблизительно равными, масса Марса в 10 раз меньше. Для низких круговых орбит в принятой модели
\[\frac{T_{\text{М}}}{T_{\text{З}}}\approx\sqrt{\frac{\rho_{\text{З}}}{\rho_{\text{М}}}}\approx 1.\]
Периоды приблизительно равны. Это результат допущений задачи, а реальные плотности планет различаются.
\clearpage
\section*{Алгоритм: решение задач на круговое движение}
\vfill
\begin{center}
\begin{tikzpicture}[node distance=9mm,>=Stealth,
flowstep/.style={draw=accent,rounded corners=3pt,align=center,text width=9.5cm,minimum height=12mm,inner sep=5pt},
term/.style={draw=darkaccent,ellipse,align=center,minimum width=3cm,minimum height=10mm},
dec/.style={diamond,draw=accent,aspect=3,align=center,inner sep=2pt,text width=4.7cm},
arr/.style={->,thick,darkaccent}]
\node[term] (start) {Начало};
\node[flowstep,below=of start] (si) {Запишите данные и искомую величину.\\ Переведите длины, время и массы в СИ.};
\node[dec,below=of si] (q) {Спутник на круговой орбите?};
\node[flowstep,text width=5.9cm,below left=12mm and -12mm of q] (grav) {$r=R_0+h$\\ $GM/r^2=v^2/r$};
\node[flowstep,text width=5.9cm,below right=12mm and -12mm of q] (kin) {Кинематика: $v=\omega r$.\\ Тяготение: $F=GMm/r^2$.\\ Для колёс: $v_1=v_2$.};
\node[flowstep,below=34mm of q] (formula) {Выберите формулу с искомой величиной.\\ При сравнении составьте отношение.};
\node[flowstep,below=of formula] (solve) {Выразите неизвестную символически,\\ затем подставьте данные (если они даны).};
\node[flowstep,below=of solve] (check) {Проверьте единицы, знак и смысл результата.\\ Запишите ответ с единицами.};
\node[term,below=of check] (end) {Конец};
\draw[arr] (start)--(si);\draw[arr] (si)--(q);
\draw[arr] (q.west)-| node[pos=.25,above]{да} (grav.north);
\draw[arr] (q.east)-| node[pos=.25,above]{нет} (kin.north);
\draw[arr] (grav.south) -- ++(0,-5mm) -| ([xshift=-23mm]formula.north);\draw[arr] (kin.south) -- ++(0,-5mm) -| ([xshift=23mm]formula.north);
\draw[arr] (formula)--(solve);\draw[arr] (solve)--(check);\draw[arr] (check)--(end);
\end{tikzpicture}
\end{center}
\vfill
\clearpage
\section{Самопроверка и типичные ошибки}
\begin{itemize}[itemsep=8pt]
\item[$\square$] Различаю период $T$, частоту $\nu$ и угловую скорость $\omega$. Знаю, что $1\unit{Гц}=1\unit{с}^{-1}$, а один оборот соответствует $2\pi$ радиан.
\item[$\square$] Помню, что $v$ направлена по касательной, а $\vec a_{\text{ц}}$ --- к центру окружности.
\item[$\square$] В задаче о шестернях сравниваю линейные скорости в точке контакта; угловые скорости могут отличаться.
\item[$\square$] При сравнении ускорений уточняю, какая величина постоянна: $v$ или $\omega$.
\item[$\square$] Для силы тяготения использую расстояние между центрами. В формуле орбиты подставляю $R_0+h$, причём в знаменателе силы стоит $(R_0+h)^2$.
\item[$\square$] Различаю гравитационную постоянную $G$ и ускорение свободного падения $g$. Ускорение на орбите обеспечивается гравитацией: $a_{\text{ц}}=g(h)$ для круговой орбиты.
\item[$\square$] При задачах на отношения сначала записываю формулы для обоих состояний и сокращаю общие величины.
\item[$\square$] При низкой орбите применяю $r\approx R_0$ только при соответствующем условии задачи.
\end{itemize}
\subsection*{Проверка размерностей}
\begin{center}
\renewcommand{\arraystretch}{1.4}
\begin{tabularx}{\linewidth}{@{}l X@{}}
\toprule Формула & Проверка в СИ\\\midrule
$v=\omega r$ & $\text{с}^{-1}\cdot\text{м}=\text{м}/\text{с}$ (радиан безразмерен)\\
$a=v^2/r$ & $(\text{м}/\text{с})^2/\text{м}=\text{м}/\text{с}^2$\\
$GM/r^2$ & $\text{Н}\cdot\text{м}^2/\text{кг}^2\cdot\text{кг}/\text{м}^2=\text{Н}/\text{кг}=\text{м}/\text{с}^2$\\
$\sqrt{GM/r}$ & $\sqrt{\text{м}^2/\text{с}^2}=\text{м}/\text{с}$\\\bottomrule
\end{tabularx}
\end{center}
\key{Перед вычислениями:} $1\unit{км}=10^3\unit{м}$, $1\unit{см}=10^{-2}\unit{м}$, $1\unit{мин}=60\unit{с}$. Округляйте результат только в конце.
\clearpage
\section*{Тренировочный блок}
Во всех задачах записывайте формулы и единицы. Для круговых орбит используйте модель неподвижной сферической планеты без сопротивления среды. При необходимости принимайте $G=6{,}7\cdot10^{-11}\unit{Н}\cdot\text{м}^2/\text{кг}^2$.
\subsection*{Базовый уровень}
\begin{enumerate}[itemsep=8pt]
\item Колесо совершает 45 оборотов за 15 с. Найдите частоту и период вращения.
\item Точка движется по окружности радиусом 40 см с периодом 2 с. Найдите угловую и линейную скорости.
\item Скорость тела, движущегося по окружности, увеличилась в 5 раз при неизменном радиусе. Во сколько раз изменилось центростремительное ускорение?
\end{enumerate}
\subsection*{Средний уровень}
\begin{enumerate}[resume,itemsep=8pt]
\item Две внешние шестерни радиусами 18 и 6 см сцеплены без проскальзывания. Большая делает 24 оборота за 12 с. Найдите частоту малой шестерни.
\item Массы двух тел увеличили в 2 и 3 раза, а расстояние между их центрами --- в 6 раз. Как изменилась сила притяжения?
\item Масса планеты B в 9 раз больше массы планеты A, её радиус в 3 раза больше. Сравните ускорения свободного падения у поверхностей.
\item Спутник движется по круговой орбите радиусом $8\cdot10^6$ м вокруг планеты массой $4{,}8\cdot10^{24}$ кг. Найдите скорость и центростремительное ускорение (до трёх значащих цифр).
\end{enumerate}
\clearpage
\subsection*{Повышенный уровень}
\begin{enumerate}[resume,itemsep=16pt]
\item У сферической планеты радиус $R_0=6000$ км, масса $M=6\cdot10^{24}$ кг. Спутник находится на круговой орбите на высоте 4000 км. Найдите радиус орбиты, угловую скорость и период (до трёх значащих цифр).
\vspace{12mm}
\item Два спутника вращаются вокруг одной планеты по круговым орбитам. Радиус второй орбиты в 9 раз больше радиуса первой. Во сколько раз отличаются линейные скорости, частоты и периоды обращения? Укажите, какая величина больше и у какого спутника.
\vspace{12mm}
\item Сферическая планета B имеет среднюю плотность в 4 раза больше плотности планеты A, а её масса в 32 раза больше. Найдите отношение радиусов $R_B/R_A$ и отношение периодов спутников $T_B/T_A$ на низких круговых орбитах, считая их радиусы равными радиусам соответствующих планет.
\end{enumerate}
\clearpage
\section*{Ответы}
\begin{center}
\renewcommand{\arraystretch}{1.65}
\begin{tabularx}{\linewidth}{@{}r X@{}}
\toprule № & Ответ\\\midrule
1 & $\nu=3\unit{Гц}$; $T=1/3\unit{с}$.\\
2 & $\omega=\pi\unit{рад/с}$; $v=0{,}4\pi\unit{м/с}$.\\
3 & Увеличилось в 25 раз.\\
4 & $6\unit{Гц}$.\\
5 & Уменьшилась в 6 раз.\\
6 & $g_B=g_A$.\\
7 & $v\approx 6{,}34\cdot10^3\unit{м/с}$; $a_{\text{ц}}\approx 5{,}03\unit{м/с}^2$.\\
8 & $r=10^7\unit{м}$; $\omega\approx 6{,}34\cdot10^{-4}\unit{рад/с}$; $T\approx9{,}91\cdot10^3\unit{с}$.\\
9 & $v_1/v_2=3$; $\nu_1/\nu_2=27$; $T_2/T_1=27$.\\
10 & $R_B/R_A=2$; $T_B/T_A=1/2$.\\\bottomrule
\end{tabularx}
\end{center}
\vspace{7mm}
\subsection*{Итоговая памятка}
\begin{center}
\renewcommand{\arraystretch}{1.35}
\begin{tabularx}{\linewidth}{@{}l X@{}}
\toprule Сигнал в задаче & Полезное действие\\\midrule
Обороты и время & Начните с $\nu=N/t$ и $T=1/\nu$.\\
Сцепленные колёса & Приравняйте $\omega_1r_1$ и $\omega_2r_2$.\\
Изменение расстояния & Используйте отношение формул.\\
Высота над планетой & Найдите $r=R_0+h$.\\
Круговая орбита & Приравняйте $GM/r^2$ и $v^2/r$.\\
Низкая орбита и плотность & Подставьте $M=4\pi R_0^3\rho/3$ в формулу периода.\\\bottomrule
\end{tabularx}
\end{center}
\end{document}
